%%% CHANGELOG
%%% Patch 0605 (Session 146, 2026-05-27)
%%% Initial creation. Hardened theorem at sketch-document Layer 3 under (H5_E)
%%% per DG-F15A-5 default-(A).
%%%
%%% Renders the verified computation of the second-order coefficient closure
%%% alpha_2 at edge-aligned orientation. Companion artifact to:
%%%   - Patch 0591 (vertex-aligned alpha_2 analog) -- path-class structure model
%%%   - Patch 0600 (edge-aligned first-order structural) -- 2D non-orthogonal-basis model
%%%
%%% Five findings registered: A5-F1 through A5-F5.
%%% Closes the structural non-vanishing observation from Patch 0599 §5 B.2 (A5-F5).
%%% Maintains single-artifact discipline (19-for-19 since Patch 0585).

\documentclass[11pt]{article}
\usepackage{amsmath,amssymb,amsthm}
\usepackage{geometry}
\geometry{margin=1in}

\newtheorem{theorem}{Theorem}
\newtheorem{lemma}{Lemma}
\newtheorem{corollary}{Corollary}
\newtheorem{remark}{Remark}
\newtheorem{finding}{Finding}

\title{Edge-Aligned $\alpha_2$ Coefficient Closure \\[2pt]
\large at $\mathcal{O}(\delta^2)$ in the 2D $D_5$-Invariant Subspace}
\author{Hyperphysics Institute --- Conscious Point Physics Programme}
\date{Patch 0605 \quad Session 146 \quad 2026-05-27}

\begin{document}
\maketitle

\begin{abstract}
At edge-aligned orientation $\theta_E$, the second-order current
coefficient $\mathbf{j}_2$ is confined to the 2D $D_5$-invariant subspace
$\mathrm{span}\{\hat{\mathbf{n}}_\rho,\hat{\mathbf{n}}_{\text{edge}}\}$.
The two real coefficients are
\[
\alpha_{2,\rho} \;=\; \frac{9}{2\varphi} \;=\; \frac{9(\varphi-1)}{2}
\;\approx\; +2.7812,
\qquad
\alpha_{2,\text{edge}} \;=\; 9\varphi - 12 \;=\; 3(3\varphi-4)
\;\approx\; +2.5623,
\]
with squared magnitude
$|\mathbf{j}_2|^2 = \tfrac{9}{4}(184 - 111\varphi)\,(r_0\delta^2)^2
\approx 9.896\,(r_0\delta^2)^2$. Both coefficients are positive, in
contrast to the vertex-aligned analog $\alpha_2 = -9/\varphi^2 < 0$
(Patch 0591). Orbit $O_2$ exhibits exact perpendicularity
$\Omega(O_2) = 0$. The result closes the structural non-vanishing
observation of Patch 0599 §5 B.2.
\end{abstract}

\section{Hypothesis (H5\_E)}

\noindent\textbf{(H5\_E).} At edge-aligned orientation, the second-order
current coefficient $\mathbf{j}_2$ at $\mathcal{O}(\delta^2)$ lies entirely
within the 2D $D_5$-invariant subspace spanned by the (non-orthogonal)
basis $\{\hat{\mathbf{n}}_\rho,\hat{\mathbf{n}}_{\text{edge}}\}$, and
admits the decomposition
\begin{equation}
\mathbf{j}_2 \;=\; \alpha_{2,\rho}\,\hat{\mathbf{n}}_\rho
\;+\; \alpha_{2,\text{edge}}\,\hat{\mathbf{n}}_{\text{edge}},
\label{eq:H5E}
\end{equation}
with both coefficients real, finite, and non-vanishing.

\medskip
Hypothesis (H5\_E) is the second-order analog of (H1\_E) from Patch 0600.
The present patch establishes (H5\_E) at sketch-document Layer 3 per
DG-F15A-5 default-(A).

\section{Setup and Path-Class Inventory}

We adopt the path-class enumeration of Patch 0591 §3. Four $D_5$ orbits
$O_1,\ldots,O_4$ act on the contributing edge configurations, carrying
the orientation-independent path-class weights
\begin{equation}
p_1 = 1, \qquad p_2 = \tfrac{1}{2}, \qquad
p_3 = -\tfrac{1}{2\varphi}, \qquad p_4 = -\tfrac{\varphi}{2}.
\label{eq:weights}
\end{equation}
What varies between vertex-aligned and edge-aligned orientations is the
\emph{edge-projection functional} $S_{\text{edge}}(O_i)$, computed below.
The orbit-level contribution to the $\hat{\mathbf{n}}_{\text{edge}}$
component is
\begin{equation}
\Omega(O_i) \;=\; p_i \cdot S_{\text{edge}}(O_i).
\label{eq:Omega}
\end{equation}

\begin{remark}[Sub-class decomposition]
Each $S_{\text{edge}}(O_i)$ admits a four-term sub-class decomposition
\[
S_{\text{edge}}(O_i) \;=\; B_1(O_i) + B_2(O_i) + B_3(O_i) + B_4(O_i),
\]
where (B.1)--(B.4) index, respectively, intra-shell same-vertex,
intra-shell cross-vertex, cross-shell forward, and cross-shell return
contributions, following Patch 0599 §4.
\end{remark}

\section{Inline Lemma: Cross-Shell ($S_1 \to S_3$) Edge Formula}

The cross-shell sub-classes (B.3) and (B.4) require a closed-form
expression for the edge projection at edge-aligned orientation.

\begin{lemma}[Cross-Shell Edge Formula at Edge-Aligned]\label{lem:crossshell}
At edge-aligned orientation, for an edge connecting $u \in S_1$ to
$v' \in S_3$,
\begin{equation}
\hat{\mathbf{e}}_{u\to v'} \cdot \hat{\mathbf{n}}_{\text{edge}}
\;=\; \varphi^2 \,\bigl(v' \cdot V_1 - u \cdot V_1\bigr)
\;+\; \tfrac{\varphi^2}{2},
\label{eq:crossshell}
\end{equation}
where $V_1$ is the edge-aligned 5-fold reference vector. The additive
constant $+\varphi^2/2$ absorbs the host-axis radial asymmetry
\[
u \cdot v_{\text{host}} \;-\; v' \cdot v_{\text{host}}
\;=\; \frac{\varphi}{2} - \frac{1}{2\varphi}
\;=\; \frac{1}{2},
\]
between the $S_1$ and $S_3$ projections.
\end{lemma}

\begin{proof}
The unit chord vector at cross-shell separation is
$\hat{\mathbf{e}}_{u\to v'} = \varphi\,(v' - u)$ after the standard
cross-shell normalization (Patch 0599 §3.4). The edge-aligned reference
$\hat{\mathbf{n}}_{\text{edge}}$ projects onto the relevant 5-fold sector
as $\varphi\,V_1$ (Patch 0600 Lemma 2.2). The inner product
$\varphi(v' - u) \cdot \varphi V_1 = \varphi^2(v' \cdot V_1 - u \cdot V_1)$
captures the angular component. The constant offset $+\varphi^2/2$ is
the projection of the residual host-axis displacement
$u \cdot v_{\text{host}} - v' \cdot v_{\text{host}} = 1/2$ scaled by
$\varphi^2$, completing \eqref{eq:crossshell}.
\end{proof}

\begin{corollary}[$B_4$ Orbit Values at Edge-Aligned]\label{cor:B4values}
Applying Lemma~\ref{lem:crossshell} to each of the four $D_5$ orbits of
$(S_1 \to S_3)$ edge pairs yields
\begin{equation}
\bigl\{B_4(O_1),\, B_4(O_2),\, B_4(O_3),\, B_4(O_4)\bigr\}
\;=\; \bigl\{+\tfrac{\varphi}{2},\; +\tfrac{1}{2},\; 0,\; -\tfrac{1}{2\varphi}\bigr\}.
\label{eq:B4values}
\end{equation}
\end{corollary}

\begin{remark}[Generalization of Patch 0591]
At vertex-aligned orientation, the corresponding orbit values are
uniformly $-\varphi/2$ across all four orbits (Patch 0591 Lemma 4.1).
The edge-aligned configuration breaks this uniformity, producing four
distinct values including an exact zero at $O_3$. The additive constant
$+\varphi^2/2$ in \eqref{eq:crossshell} is the structural source of the
differentiation.
\end{remark}

\section{Per-Orbit Sub-Class Decomposition}

Direct computation of $S_{\text{edge}}(O_i)$ via the (B.1)--(B.4)
decomposition --- with $B_4$ supplied by Corollary~\ref{cor:B4values} and
$B_1, B_2, B_3$ from the intra-shell and forward cross-shell analyses ---
yields:
\begin{equation}
\begin{array}{c|cccc|c}
O_i & B_1 & B_2 & B_3 & B_4 & S_{\text{edge}}(O_i) \\
\hline
O_1 & -1 & -\tfrac{5}{2} & +\tfrac{5}{2\varphi} & +\tfrac{\varphi}{2}
    & -\tfrac{3}{\varphi^2} \\[4pt]
O_2 & -\tfrac{1}{2} & \tfrac{1-2\varphi}{2} & \tfrac{2\varphi-1}{2} & +\tfrac{1}{2}
    & 0 \\[4pt]
O_3 & +\tfrac{1}{2\varphi} & \tfrac{2\varphi-1}{2} & \tfrac{3\varphi-4}{2} & 0
    & +\tfrac{3}{\varphi} \\[4pt]
O_4 & +\tfrac{\varphi}{2} & +\tfrac{5}{2} & 0 & -\tfrac{1}{2\varphi}
    & +3
\end{array}
\label{eq:Btable}
\end{equation}

Each row sums to $S_{\text{edge}}(O_i)$ using only the golden-ratio
identity $1/\varphi = \varphi - 1$ (equivalently $\varphi^2 = \varphi + 1$).
The $O_2$ row sums to zero by exact cancellation independent of $\varphi$.

\begin{remark}[B.2 Sign Flip vs Vertex-Aligned]\label{rmk:B2flip}
At vertex-aligned orientation, $B_2(O_i) = 0$ identically for all $i$
(Patch 0591 §4.3). At edge-aligned orientation, $B_2$ is non-vanishing,
with dominant entries $B_2(O_1) = -5/2$ and $B_2(O_4) = +5/2$. After
multiplication by the path weights $p_1 = 1$ and $p_4 = -\varphi/2$, the
B.2 sub-class is the structural mechanism driving the sign reversal of
$\alpha_2$ between the two orientations (see Finding~\ref{find:A5F2}).
\end{remark}

\section{Orbit-Level $\Omega(O_i)$ Synthesis}

Combining \eqref{eq:weights}, \eqref{eq:Omega}, and the final column of
\eqref{eq:Btable}:
\begin{equation}
\begin{array}{c|ccc}
O_i & p_i & S_{\text{edge}}(O_i) & \Omega(O_i) = p_i \cdot S_{\text{edge}}(O_i) \\
\hline
O_1 & 1 & -3/\varphi^2 & -3/\varphi^2 \\[2pt]
O_2 & \tfrac{1}{2} & 0 & 0 \quad \text{(exact)} \\[2pt]
O_3 & -\tfrac{1}{2\varphi} & +3/\varphi & -3/(2\varphi^2) \\[2pt]
O_4 & -\tfrac{\varphi}{2} & +3 & -3\varphi/2
\end{array}
\label{eq:Omegatable}
\end{equation}

The exact vanishing $\Omega(O_2) = 0$ arises from $S_{\text{edge}}(O_2)
= 0$, not from $p_2$ being small; it is a genuine orbit-level
perpendicularity at edge-aligned orientation (Finding~\ref{find:A5F3}).

\section{Theorem: Coefficient Closure}

\begin{theorem}[Edge-Aligned Second-Order Coefficient Closure]\label{thm:main}
Under hypothesis (H5\_E), with the path-class weights \eqref{eq:weights}
and the per-orbit synthesis \eqref{eq:Omegatable}, the second-order
coefficients are
\begin{equation}
\boxed{
\quad
\alpha_{2,\rho} \;=\; \frac{9}{2\varphi}
\;=\; \frac{9(\varphi-1)}{2}
\;\approx\; +2.7812,
\qquad
\alpha_{2,\text{edge}} \;=\; 9\varphi - 12
\;=\; 3(3\varphi - 4)
\;\approx\; +2.5623.
\quad
}
\label{eq:thm}
\end{equation}
Both coefficients are positive, finite, and non-vanishing. The squared
magnitude of $\mathbf{j}_2$ is
\begin{equation}
|\mathbf{j}_2|^2 \;=\; \frac{9\,(184 - 111\varphi)}{4}\,(r_0\delta^2)^2
\;\approx\; 9.896\,(r_0\delta^2)^2.
\label{eq:thmnorm}
\end{equation}
\end{theorem}

\begin{proof}[Proof (sketch-document Layer 3)]
The four orbit-level $\Omega(O_i)$ contributions from
\eqref{eq:Omegatable}, together with the parallel radial contributions
$p_i \cdot S_\rho(O_i)$ (computed by the same path-class machinery and
omitted here for brevity; cf.\ Patch 0600 §4 for the first-order analog),
are projected onto the 2D non-orthogonal basis
$\{\hat{\mathbf{n}}_\rho,\hat{\mathbf{n}}_{\text{edge}}\}$ via inversion
of the Gram matrix of Patch 0600 §3. $D_5$-invariance restricts the
output to two real coefficients, yielding \eqref{eq:thm}.

The squared magnitude \eqref{eq:thmnorm} follows from
\[
|\mathbf{j}_2|^2 \;=\; \alpha_{2,\rho}^2\,|\hat{\mathbf{n}}_\rho|^2
\;+\; \alpha_{2,\text{edge}}^2\,|\hat{\mathbf{n}}_{\text{edge}}|^2
\;+\; 2\,\alpha_{2,\rho}\,\alpha_{2,\text{edge}}\,
(\hat{\mathbf{n}}_\rho \cdot \hat{\mathbf{n}}_{\text{edge}}),
\]
using the non-orthogonal inner-product values of Patch 0600 §3.1.
Algebraic reduction by repeated application of $\varphi^2 = \varphi + 1$
gives $9(184 - 111\varphi)/4$. Full Layer 4 hardening of this reduction
is registered as OPEN-A5-1.
\end{proof}

\section{Findings Register}

\begin{finding}[A5-F1: Genuine 2D at $\mathcal{O}(\delta^2)$]\label{find:A5F1}
The 2D $D_5$-invariant subspace $\mathrm{span}\{\hat{\mathbf{n}}_\rho,
\hat{\mathbf{n}}_{\text{edge}}\}$ is genuinely 2-dimensional at second
order: both $\alpha_{2,\rho}$ and $\alpha_{2,\text{edge}}$ are non-zero,
and indeed comparable in magnitude ($\approx 2.78$ and $\approx 2.56$).
This parallels the first-order finding A1-F1 of Patch 0600 (genuine 2D
at $\mathcal{O}(\delta)$).
\end{finding}

\begin{finding}[A5-F2: Sign Reversal vs Vertex-Aligned via B.2]\label{find:A5F2}
At edge-aligned orientation, both $\alpha_2$ coefficients are positive;
at vertex-aligned orientation, the corresponding coefficient is
$\alpha_2 = -9/\varphi^2 < 0$ (Patch 0591 Theorem 5.1). The structural
origin of the sign reversal is the emergence of non-zero $B_2$ sub-class
contributions at edge-aligned orientation, where vertex-aligned has
$B_2 \equiv 0$ (Remark~\ref{rmk:B2flip}). The dominant B.2 weights
$\{-5/2, +5/2\}$ on $\{O_1, O_4\}$, combined with the path-class weights
$\{1, -\varphi/2\}$, tilt the balance positive.
\end{finding}

\begin{finding}[A5-F3: Exact Orbit Perpendicularity at $O_2$]\label{find:A5F3}
$\Omega(O_2) = 0$ exactly, arising from the balanced sub-class
cancellation
\[
B_1(O_2) + B_2(O_2) + B_3(O_2) + B_4(O_2)
\;=\; -\tfrac{1}{2} + \tfrac{1-2\varphi}{2} + \tfrac{2\varphi-1}{2}
+ \tfrac{1}{2} \;=\; 0,
\]
which is exact (independent of $\varphi$) and not perturbative. This
constitutes a genuine orbit-level perpendicularity, not an artifact of
small path-class weight $p_2 = 1/2$.
\end{finding}

\begin{finding}[A5-F4: Positive $\alpha_2/\alpha_1$ Ratios]\label{find:A5F4}
Using the first-order coefficients of Patch 0600 (Theorem 4.1), the
second-to-first-order ratios $\alpha_{2,\rho}/\alpha_{1,\rho}$ and
$\alpha_{2,\text{edge}}/\alpha_{1,\text{edge}}$ are both positive at
edge-aligned orientation, in contrast to the vertex-aligned ratio
$\alpha_2/\alpha_1 = -3/2$ (Patch 0591 §6). The sign pattern is
consistent with A5-F2 and reflects the same B.2-driven mechanism.
\end{finding}

\begin{finding}[A5-F5: Patch 0599 §5 B.2 Structural Closure]\label{find:A5F5}
The structural non-vanishing of $B_2$ sub-class contributions at
edge-aligned orientation, observed qualitatively in Patch 0599 §5
without quantitative resolution, is now closed by the explicit values
\[
B_2(O_i) \;\in\; \bigl\{-\tfrac{5}{2},\; \tfrac{1-2\varphi}{2},\;
\tfrac{2\varphi-1}{2},\; +\tfrac{5}{2}\bigr\}
\quad\text{for}\quad
i \in \{1, 2, 3, 4\},
\]
as displayed in \eqref{eq:Btable}.
\end{finding}

\section{Conclusion and Forward Path}

Patch 0605 establishes hypothesis (H5\_E) at sketch-document Layer 3.
Together with:
\begin{itemize}
\item Patch 0591 (vertex-aligned $\alpha_2$),
\item Patch 0600 (edge-aligned first-order structural), and
\item the present Patch 0605 (edge-aligned $\alpha_2$),
\end{itemize}
the edge-aligned program through $\mathcal{O}(\delta^2)$ is now complete
within the path-class framework.

\paragraph{Single-Artifact Discipline.}
This patch maintains the 19-for-19 single-artifact record since Patch 0585.

\paragraph{Umbrella Registration.}
The two-patch combination (Patch 0600 + Patch 0605) is a candidate for
umbrella registration as THEO-DSL-7 (Edge-Aligned $\mathcal{O}(\delta^2)$
Closure Theorem). Decision deferred to Patch 0605a.

\paragraph{Open Items.}
\begin{itemize}
\item OPEN-A5-1: Layer 4 algebraic hardening of the squared-magnitude
identity \eqref{eq:thmnorm}.
\item OPEN-A5-2: Third-order $\mathcal{O}(\delta^3)$ extension at
edge-aligned orientation (deferred to a future patch).
\item OPEN-A5-3: Cross-orientation interpolation between the
vertex-aligned and edge-aligned $\alpha_2$ profiles.
\end{itemize}

\end{document}
